\documentclass[11pt]{article}
\usepackage[a4paper,margin=27mm]{geometry}
\usepackage{amsmath,amssymb,mathtools}
\usepackage{fontspec,unicode-math}
\setmainfont{Latin Modern Roman}
\setmathfont{Latin Modern Math}
\usepackage{xurl,hyperref}
\hypersetup{colorlinks=true,urlcolor=blue,linkcolor=blue}
\newcommand{\sref}[2]{\hyperlink{source:#1:#2}{[S#2]}}
\newcommand{\cataloguescope}[1]{\paragraph{Mathematical question.} #1}
\begin{document}
% UnsolvedMath ID 8700006; source catalogue code AMR-086-0006
\setcounter{section}{8700005}
\section{Conjecture 1.7}\label{8700006}
{\small\sffamily UnsolvedMath ID 8700006 · Graph Theory}
\subsection{Definitions and mathematical statement}

\paragraph{Shared notation.}$\mathbb N=\{1,2,\ldots\}$, and $\mathbb Z,\mathbb Q,\mathbb R,\mathbb C$ denote the integers, rationals, reals and complex numbers. Any other conventions are specified locally.


For a positive integer $N$, let $\omega(N)$ be the number of its distinct prime divisors. Consider $n,k\in\mathbb N$, so the consecutive block $n+1,\ldots,n+k$ lies above one.
\paragraph{Statement recorded in the supplied source.}
If the interval $[n+1,n+k]$ contains no prime, must \[\omega\bigl((n+1)\cdots(n+k)\bigr)\ge k?\] \sref{8700006}{2}
\subsection{Short English statement}
Must a consecutive block of composite numbers have at least as many different prime factors collectively as it has terms?
\subsection{Sources}
\begin{description}
\item[\hypertarget{source:8700006:1}{S1}] ULAM.AI and the UnsolvedMath Contributors, UnsolvedMath: A Curated Collection of Open Mathematics Problems, record 8700006 (AMR-086-0006). CC BY 4.0. \url{https://huggingface.co/datasets/ulamai/UnsolvedMath}.
\item[\hypertarget{source:8700006:2}{S2}] Michel Waldschmidt, Open Diophantine Problems, arXiv:math/0312440v2 (24 January 2004), Conjecture 1.7. \url{https://arxiv.org/abs/math/0312440}.
\end{description}
\label{end:8700006}
\end{document}
